% answers_short.tex —— 2016 年·山东卷·理科 · 简短答案速查
% 由 tools/make_short_answers.py 自 answers.tex 的 \gksolhead/\gkanswer 提取，
% 并与 2016/2016_数学真题_山东理_答案要点.md 逐题比对后生成，勿手改（改源 MD 后重跑）。

\begin{gknotice}
\textbf{使用说明}：本节只给最终答案便于快速核对，完整推导见解析版。答案为编写者独立推导，\textbf{非官方标准答案}。
\end{gknotice}

\gkpart{选择题}
\begin{center}\begin{tabular}{ccccccccccc}
\toprule
题号 & 1 & 2 & 3 & 4 & 5 & 6 & 7 & 8 & 9 & 10 \\
\midrule
答案 & B & C & D & C & C & A & B & B & D & A \\
\bottomrule
\end{tabular}\end{center}

\gkpart{填空题}
\begin{itemize}[leftmargin=2.6em]
  \item[\textbf{11.}] \(3\)
  \item[\textbf{12.}] \(-2\)
  \item[\textbf{13.}] \(2\)
  \item[\textbf{14.}] \(\dfrac{3}{4}\)
  \item[\textbf{15.}] \(m\in(3,+\infty)\)
\end{itemize}

\gkpart{解答题}
\begin{itemize}[leftmargin=2.6em]
  \item[\textbf{16.}] (1) 见解析版；(2) \(\cos C\) 的最小值为 \(\dfrac{1}{2}\)（当且仅当 \(a=b\) 时取得）
  \item[\textbf{17.}] (1) 见解析版；(2) 二面角 \(F-BC-A\) 的余弦值为 \(\dfrac{\sqrt{7}}{7}\)
  \item[\textbf{18.}] (1) \(b_{n}=3n+1\)；(2) \(T_{n}=3n\cdot 2^{n+2}\)
  \item[\textbf{19.}] (1) 两轮至少猜对 \(3\) 个成语的概率为 \(\dfrac{2}{3}\)；(2) \(P(X=0)=\dfrac{1}{144}\)，\(P(X=1)=\dfrac{5}{72}\)，\(P(X=2)=\dfrac{25}{144}\)，\(P(X=3)=\dfrac{1}{12}\)，\(P(X=4)=\dfrac{5}{12}\)，\(P(X=6)=\dfrac{1}{4}\)；\(EX=\dfrac{23}{6}\)
  \item[\textbf{20.}] (1) \(f'(x)=\dfrac{(x-1)(ax^{2}-2)}{x^{3}}\)（\(x>0\)）。当 \(a\leqslant 0\) 时，\(f(x)\) 在 \((0,1)\) 上递增、在 \((1,+\infty)\) 上递减；当 \(0<a<2\) 时，在 \((0,1)\) 与 \(\left(\sqrt{\dfrac{2}{a}},+\infty\right)\) 上递增、在 \(\left(1,\sqrt{\dfrac{2}{a}}\right)\) 上递减；当 \(a=2\) 时，在 \((0,+\infty)\) 上递增；当 \(a>2\) 时，在 \(\left(0,\sqrt{\dfrac{2}{a}}\right)\) 上递减、在 \(\left(\sqrt{\dfrac{2}{a}},+\infty\right)\) 上递增；(2) 见解析版
  \item[\textbf{21.}] (1) \(\dfrac{x^{2}}{1}+\dfrac{y^{2}}{\frac{1}{4}}=1\)（即 \(x^{2}+4y^{2}=1\)）；(2) ① 点 \(M\) 在定直线 \(y=-\dfrac{1}{4}\) 上（证明见解析版）；② \(\dfrac{S_{1}}{S_{2}}\) 的最大值为 \(\dfrac{9}{4}\)，取得最大值时 \(P\left(\dfrac{\sqrt{2}}{2},\dfrac{1}{4}\right)\)
\end{itemize}
